\setcounter{section}{6}

\Needspace{5\baselineskip}
\hypertarget{ux89c6ux9891ux751fux6210ux6a21ux578bux7684ux4e0aux4e0bux6587ux5de5ux7a0b}{%
\section{视频生成模型的上下文工程}\label{ux89c6ux9891ux751fux6210ux6a21ux578bux7684ux4e0aux4e0bux6587ux5de5ux7a0b}}

\Needspace{5\baselineskip}
\hypertarget{ux4e00ux4e0aux4e0bux6587ux7684ux65f6ux7a7aux538bux7f29}{%
\subsection{一、上下文的时空压缩}\label{ux4e00ux4e0aux4e0bux6587ux7684ux65f6ux7a7aux538bux7f29}}

\Needspace{5\baselineskip}
\hypertarget{ux4e00ux95eeux9898ux80ccux666f-5}{%
\subsubsection{（一）问题背景}\label{ux4e00ux95eeux9898ux80ccux666f-5}}

\textbf{要解决的问题}：在自回归生成长视频时，随着历史帧 \(z_c\) 数量的不断增加，Transformer（DiT）的 Context Length 会呈线性甚至平方级增长，导致显存爆炸和推理速度断崖式下降。传统的滑动窗口（Sliding Window）会丢失长程历史信息，而直接丢弃历史帧则会导致场景不一致。

\Needspace{5\baselineskip}
\hypertarget{ux4e8cux65f6ux95f4ux4e0eux7a7aux95f4ux538bux7f29}{%
\subsubsection{（二）时间与空间压缩}\label{ux4e8cux65f6ux95f4ux4e0eux7a7aux95f4ux538bux7f29}}

\Needspace{5\baselineskip}
对于历史帧 \(z_c\)，模型首先在时间维度上进行每 32 帧抽样 1 帧的随机采样，然后使用高压缩比的 Patchify 操作。距离当前预测帧越远的历史帧，下采样率越大：

\begin{itemize}
\tightlist
\item
  距离 \(t-1\) 到 \(t-2\) 的帧：下采样率为 \((1,2,2)\)。
\item
  距离 \(t-3\) 到 \(t-6\) 的帧：下采样率为 \((1,4,4)\)。
\item
  距离 \(t-7\) 到 \(t-23\) 的帧：下采样率为 \((1,8,8)\)。
\end{itemize}

由于世界在时空上的近似连续性特质，这种下采样对世界模型性能的影响较小，但能大大减轻显存压力。

\Needspace{5\baselineskip}
\hypertarget{ux4e09ux901aux9053ux538bux7f29ux4e0eux7ebfux6027ux6ce8ux610fux529b}{%
\subsubsection{（三）通道压缩与线性注意力}\label{ux4e09ux901aux9053ux538bux7f29ux4e0eux7ebfux6027ux6ce8ux610fux529b}}

进一步地，历史帧 \(z_c\) 经过压缩率为 \((8,4,4)\) 的 Patchify，并将其通道维度压缩至 96，得到 \(z_{linear}\)。

在 DiT Block 内部，当前视频 Token \(z^l\) 经过交叉注意力层后，提取出预测帧 \(z_p^l\)，并与 \(z_{linear}\) 拼接，得到融合特征 \(z_{fus}\)。

为了避免标准注意力机制中序列长度带来的二次方计算复杂度，模型在此处使用了\textbf{线性注意力（Linear Attention）}。

\Needspace{5\baselineskip}
\hypertarget{ux4e8ckv-cacheux4e2dux7684ux91cdux951aux70b9}{%
\subsection{二、KV Cache中的重锚点}\label{ux4e8ckv-cacheux4e2dux7684ux91cdux951aux70b9}}

自回归模型有一个通病：误差累积。生成到第100块时，它参考的是第99块和第98块；如果第99块有一点点误差，第100块就会照着变形的画面继续生成，最后导致严重漂移。

\begin{enumerate}
\def\labelenumi{\arabic{enumi}.}
\item
  \textbf{注意力缓存（KV Cache）的常规逻辑}：模型在生成新帧时，会``回顾''之前的帧（即查阅 KV 缓存）。通常，模型倾向于关注最近的几帧以保持动作连贯。
\item
  \textbf{AHIS 的核心思想：``厚古薄今''}。

  \begin{itemize}
  \tightlist
  \item
    论文设定了一个容量为 \textbf{5 个 Block（区块）} 的注意力窗口。
  \item
    \textbf{身份锚点（Sink Tokens）}：将最开始生成的、质量最高、身份最完美的前 3 个 Block 永久性地锁定在缓存里（不随时间滑动而被淘汰）。这些就像是``身份证件照''。
  \item
    \textbf{滚动缓存（Rolling Tokens）}：剩下的 2 个 Block 用来存储最近刚刚生成的画面，并且随着时间不断滑动更新，用于保持当前嘴唇和动作的物理连贯性。
  \end{itemize}
\item
  \textbf{为什么叫 Anchor-Heavy（重锚点）？}

  \begin{itemize}
  \tightlist
  \item
    传统的注意力下沉（Attention Sinks，多用于大语言模型）通常只保留极少量的 Token（比如第一句话的首个词）用来稳定注意力计算。
  \item
    而这篇论文故意给身份锚点分配了极其不成比例的巨大权重（3 个 Block 对比 2 个 Block）。
  \item
    \textbf{效果}：这强迫模型在生成每一帧时，绝大部分的注意力都死死盯住最初始那张没有畸变的``完美脸（高保真表征）''，从而在根本上压制了模型去模仿最近帧里那些细微的扭曲误差。这使得数字人即使连续说话数分钟，长相也不会发生漂移。
  \end{itemize}
\end{enumerate}

\Needspace{5\baselineskip}
\hypertarget{ux4e09ux8badux7ec3ux4e0aux4e0bux6587ux6a21ux62dfux6ce8ux5165ux8befux5dee}{%
\subsection{三、训练上下文模拟注入误差}\label{ux4e09ux8badux7ec3ux4e0aux4e0bux6587ux6a21ux62dfux6ce8ux5165ux8befux5dee}}

在推理时，模型的上下文是自身已经生成的帧，难免有误差。如果训练时模型只学会了基于真实帧进行生成，则在推理时在出现误差时就容易进入OOD区域。为解决该问题，昆仑万维的Matrix-Game 3.0在训练上下文中就模拟注入误差，从而增强模型对真实场景下上下文误差的鲁棒性。

\begin{enumerate}
\def\labelenumi{\arabic{enumi}.}
\item
  \textbf{分组输入}：将视频潜变量序列分为前 \(k\) 帧（作为历史条件）和后 \(N-k\) 帧（作为当前需要预测的加噪目标）。
\item
\Needspace{5\baselineskip}
  \textbf{误差收集}：计算模型预测的干净估计值 \(\hat{x}^t\) 与真实潜变量 \(x^t\) 之间的残差：
\end{enumerate}

\[
\delta = \hat{x}^t - x^t
\]

\Needspace{4\baselineskip}
\begin{enumerate}
\def\labelenumi{\arabic{enumi}.}
\setcounter{enumi}{2}
\tightlist
\item
\Needspace{5\baselineskip}
  \textbf{误差注入}：将残差存入误差缓冲区 \(\mathcal{E}\)，在训练时按照均匀分布采样标量 \(\gamma\)，对历史潜变量进行扰动：
\end{enumerate}

\[
\tilde{x}^t = x^t + \gamma\delta
\]

\begin{enumerate}
\def\labelenumi{\arabic{enumi}.}
\setcounter{enumi}{3}
\tightlist
\item
\Needspace{5\baselineskip}
  \textbf{条件注入与优化}：通过交叉注意力（Cross-Attention）准确注入离散的键盘动作，通过自注意力（Self-Attention）注入连续的鼠标控制信号。最终的流匹配（Flow-matching）优化目标为：
\end{enumerate}

\[
\mathcal{L} = \mathbb{E}_{x,t,\epsilon,\delta}\left[\left\|\left(\epsilon - x^{k+1:N}\right) - v_{\theta}\left(x_t^{k+1:N}, t\mid \tilde{x}^{1:k}, c\right)\right\|_2^2\right]
\]

\Needspace{5\baselineskip}
\hypertarget{ux56dbux76f8ux673aux611fux77e5ux7684ux957fux7a0bux8bb0ux5fc6ux673aux5236}{%
\subsection{四、相机感知的长程记忆机制}\label{ux56dbux76f8ux673aux611fux77e5ux7684ux957fux7a0bux8bb0ux5fc6ux673aux5236}}

Matrix-Game 3.0通过设计长程记忆模块，提升了在自回归的长程生成中的空间一致性与稳定性。

为了保证玩家在虚拟世界中长时间漫游后``回头看''时，场景布局和细节不发生改变，模型引入了显式的长程记忆。

\begin{itemize}
\tightlist
\item
  \textbf{统一自注意力机制}：检索到的记忆潜变量、近期过去帧和当前加噪帧被置于同一个注意力空间中，由同一个 DiT 进行联合时空建模，避免了外部记忆分支带来的特征不对齐和收敛缓慢问题。
\item
  \textbf{相机感知检索与编码}：并非所有历史都有用。模型根据相机的位姿和视场角重叠度来检索相关的记忆帧，并使用相对 Plücker 坐标（Plücker encoding）进行显式的几何编码，帮助模型准确对齐不同视角下的同一场景。记忆通路同样采用了上述的误差注入工作流，以减少训练和推理间的偏差。
\item
\Needspace{5\baselineskip}
  \textbf{改进的旋转位置编码（RoPE）}：为了避免相隔很远的记忆帧和当前帧出现周期性的位置对齐（Positional aliasing），模型对不同注意力头引入了独立的扰动系数 \(\epsilon_h\)：
\end{itemize}

\[
\hat{\theta}_h = \theta_{base}(1 + \sigma_{\theta}\epsilon_h)
\]

\FloatBarrier
\Needspace{5\baselineskip}
\hypertarget{ux53c2ux8003ux6587ux732e-144}{%
\subsection{参考文献}\label{ux53c2ux8003ux6587ux732e-144}}

\vspace{0.25em}
\begin{itemize}
\tightlist
\item
  Wang, Z., Liu, Z., Li, J., Huang, K., Xu, B., Kang, F., An, M., Wang, P., Jiang, B., Wei, Y., Xietian, Y., Pei, J., Hu, L., Jiang, B., Xue, H., Wang, Z., Sun, H., Li, W., Ouyang, W., He, X., Liu, Y., Li, Y., \& Zhou, Y. (2026). \href{https://arxiv.org/abs/2604.08995}{Matrix-Game 3.0: Real-Time and Streaming Interactive World Model with Long-Horizon Memory}. arXiv:2604.08995.
\item
  Chen, B., Monso, D. M., Du, Y., Simchowitz, M., Tedrake, R., \& Sitzmann, V. (2024). \href{https://arxiv.org/abs/2407.01392}{Diffusion Forcing: Next-token Prediction Meets Full-Sequence Diffusion}. NeurIPS.
\item
  Mao, X., Li, Z., Li, C., et al.~(2025). \href{https://arxiv.org/abs/2512.22096}{Yume-1.5: A Text-Controlled Interactive World Generation Model}. arXiv:2512.22096.
\item
  Li, W., Pan, W., Luan, P.-C., Gao, Y., \& Alahi, A. (2025). \href{https://arxiv.org/abs/2510.09212}{Stable Video Infinity: Infinite-Length Video Generation with Error Recycling}. ICLR.
\end{itemize}

\clearpage
